\section{Introduction}

Crop-yield estimates at the district level are used in India for procurement planning, crop
insurance and early warning of agricultural stress. Weather strongly influences yield, and
machine-learning models that learn the relation between seasonal weather and yield have been
reported to predict yield with useful accuracy \cite{vanklompenburg2020,jeong2016,khaki2019}.
Random forests and gradient-boosted trees are now common choices for this task, and they have
been applied to Indian district statistics combined with gridded weather from NASA POWER
\cite{alsaber2025,manogna2026}.

Two conditions that arise when such a model is used in practice are, however, rarely tested.
First, weather inputs are often incomplete. Station outages, processing failures and errors
when sources are merged leave gaps in the daily record, and a model trained on complete inputs
must then work with fewer or distorted features. Imputation methods ranging from simple medians
to iterative model-based estimates are available \cite{little2019,troyanskaya2001,vanbuuren2011},
but in crop-yield studies they have been compared only on the gaps that a dataset happened to
contain \cite{saini2024,calin2023}. Second, a model is often applied to districts that were not
represented when it was trained. Districts differ in rainfall regime, soil and irrigation, and
when observations from the same district appear in both the training and the test set, the
reported accuracy describes familiar districts only. Spatially structured data are known to make
random cross-validation over-optimistic \cite{roberts2017,meyer2019,ploton2020}, and recent work on
United States counties and African countries has shown that yield models can lose most of their
skill on unseen regions \cite{chakravarty2026,adjei2026}.

These two problems have been studied separately, and neither has been studied for Indian
district data in combination with the other. A model may tolerate missing inputs in districts
it knows and yet fail when the inputs are incomplete \emph{and} the district is new. It also
remains unclear which imputation strategy should be preferred under each kind of gap, and how
much the random-split evaluation used in most earlier studies overstates the accuracy that can
be expected in a realistic deployment.

In the present study, these questions are addressed through a controlled experiment on 32
districts of Gujarat. Official district statistics for groundnut, bajra (pearl millet) and wheat
from 1997--98 to 2022--23 \cite{desapy} were linked with daily NASA POWER weather \cite{nasapower},
and seasonal weather features were derived for each crop's growing season. Four models were
compared under three evaluation designs (a random split, a temporal split and a district
holdout), and missing weather was simulated under three mechanisms at rates of 10 to 40\%:
random masking of features, removal of whole variable groups and contiguous outages in the daily
record. Four ways of handling the gaps were compared, and their effect was measured in familiar
and in unseen districts so that the interaction of the two problems could be estimated.
Feature selection and location features were also evaluated.

The following research questions are answered: (RQ1) how strongly missing weather degrades
prediction, and which strategy best preserves it; (RQ2) how much accuracy is lost on districts
excluded from training; (RQ3) whether missing inputs and geographic shift together cause a larger
loss than the sum of their separate effects; and (RQ4) whether feature selection improves
robustness. The code, the processed data and all predictions have been made publicly
available.\footnote{\url{https://github.com/ayaan-goel/ml-paper}}

The remainder of the paper is organised as follows. Related studies are reviewed and the research
gap is identified in Section~\ref{sec:lit}. The research design, data, preprocessing, missingness
model and training configuration are described in Section~\ref{sec:method}, and the performance
metrics are defined in Section~\ref{sec:metrics}. The results are presented in
Section~\ref{sec:results} and compared with earlier findings in Section~\ref{sec:comp}. The paper
is concluded in Section~\ref{sec:concl}.
